\lecture{Lecture 9: Morphological Operations and Thresholding}

\qa{In morphological operations, how does Erosion differ from Dilation in terms of their physical effect on foreground shapes and component connectivity?}{
\begin{itemize}
    \item \textbf{Erosion} shrinks or thins foreground objects. It is used to separate weakly connected components, as a pixel is kept only if the structuring element is \textit{fully included} in the object set ($B_z \subseteq A$).
    \item \textbf{Dilation} expands or thickens foreground objects. It is used to merge close and unconnected components, as a pixel is kept if the structuring element has \textit{at least one pixel overlapping} with the object set ($B_z \cap A \neq \emptyset$).
\end{itemize}
}

\qa{Structurally and mathematically, what is the exact difference between an Opening operation versus a Closing operation?}{
\begin{itemize}
    \item \textbf{Opening} is defined mathematically as an Erosion followed by a Dilation: $A \circ B = (A \ominus B) \oplus B$. Structurally, it smooths the contour of an object, breaks narrow isthmuses, and eliminates thin protrusions.
    \item \textbf{Closing} is defined mathematically as a Dilation followed by an Erosion: $A \bullet B = (A \oplus B) \ominus B$. Structurally, it also smooths contours but does so by fusing narrow breaks and long thin gulfs, eliminating small holes, and filling gaps.
\end{itemize}
}

\qa{When dealing with imperfect binary images, how does Opening perform versus Closing in terms of noise removal and shape repair?}{
\begin{itemize}
    \item \textbf{Opening} is best for ``opening spaces'' in the background: it removes small, bright (white) details or ``salt'' noise that are smaller than the structuring element, smoothing contours without reducing the main object's size.
    \item \textbf{Closing} is best for ``closing gaps'' in the foreground: it fills in small, dark (black) holes, narrow breaks, or ``pepper'' noise within the white objects, repairing the foreground without increasing the main object's overall size.
\end{itemize}
}

\qa{Since morphological operations are not cumulative, when is it appropriate to use [Opening then Closing] versus [Closing then Opening]?}{
\begin{itemize}
    \item Use \textbf{Opening then Closing} when the background noise is severe but the main object is relatively sturdy (e.g., a noisy fingerprint). The Opening clears the background ``salt'' noise first so it doesn't get fused to the object, then the Closing patches cracks in the ridges.
    \item Use \textbf{Closing then Opening} when the main object is highly fragmented or full of holes (e.g., poorly printed dot-matrix text). You must Close first to fuse the broken pieces together; if you Opened first, the initial erosion step would completely erase the fragile object.
\end{itemize}
}

\qa{In the context of morphological descriptions, what is the conceptual difference between the image set versus the structuring element?}{
\begin{itemize}
    \item The \textbf{image set} is the mathematical representation of the image itself. In the simplest case of a binary image, the image is defined as the set of coordinates containing the foreground pixels. Morphological operators act on this set by adding or removing pixels to/from it.
    \item The \textbf{structuring element} is the predefined ``probe'' or shape applied to the image. Its specific shape determines the direction(s) and extent to which pixels are added to or removed from the image set.
\end{itemize}
}

\qa{When segmenting an image, why might variable (local) thresholding be preferred over global thresholding?}{
\begin{itemize}
    \item \textbf{Global thresholding} applies one single threshold value to the entire image. It works well only if the image has consistent lighting and a uniform background. It fails completely when the image suffers from uneven illumination, shadows, or gradients, because no single threshold can separate the foreground from the background everywhere.
    \item \textbf{Variable (or local) thresholding} divides the image into smaller sub-regions and calculates a distinct threshold for each specific area based on its local neighborhood. This allows the algorithm to dynamically adapt to local lighting changes, successfully segmenting objects even when the background intensity drifts across the image.
\end{itemize}
}

\qa{How does Otsu's method use inter-class variance versus intra-class variance to find the optimal threshold?}{
Otsu's method evaluates every possible threshold $k$, computes class probabilities and means, and picks the $k$ that maximizes the inter-class variance $\sigma_B^2$. Since the global variance is fixed ($\sigma_G^2 = \sigma_{in}^2 + \sigma_B^2$), maximizing $\sigma_B^2$ is exactly equivalent to minimizing the intra-class variance $\sigma_{in}^2$, giving the optimal separation between the two classes.
}

\qa{In the context of Otsu's algorithm, what is the conceptual difference between the cumulative mean \texorpdfstring{$m(k)$}{m(k)} and the mean intensity value \texorpdfstring{$m_1(k)$}{m1(k)} for Class 1?}{
\begin{itemize}
    \item The \textbf{cumulative mean} $m(k) = \sum_{i=0}^{k} i p_i$ calculates the mean up to intensity level $k$, but it is normalized over the \textit{whole} image.
    \item The \textbf{mean intensity value} for Class 1, $m_1(k) = \frac{1}{P_1(k)} \sum_{i=0}^{k} i p_i$, takes that cumulative mean and divides it by $P_1(k)$. This normalizes the value specifically to the number of pixels actually belonging to Class 1, giving the true average intensity of that segment.
\end{itemize}
}

\qa{How does image noise and object size affect the standard Otsu's method versus an approach that incorporates image smoothing?}{
\begin{itemize}
    \item \textbf{Basic Otsu's method} struggles when dealing with noisy images or when target regions are very small (meaning they have a very limited impact on the histogram). Noise blurs the separation between histogram peaks.
    \item \textbf{Otsu's method with Smoothing} applies a spatial smoothing filter before calculating the histogram. This reduces noise, creates deeper valleys and clearer peaks in the histogram, and allows the algorithm to successfully segment smaller or noisier patterns.
\end{itemize}
}
